\thispagestyle{front}
\vspace*{0.5cm}
\begin{center}
  \textbf{\Large ABSTRACT}
  \addcontentsline{toc}{chapter}{Abstract}
\end{center}
\vspace{1cm}

\noindent
Most navigation systems today update routes only after a disruption has already
been detected by sensors -- they are reactive by design. This project builds an
\textbf{anticipatory} alternative for Kolkata, India.

\vspace{0.5cm}
\noindent
The system, \textbf{traffic-ai}, pulls unstructured text from nine live sources
(news RSS feeds, TomTom traffic incidents, official advisories, weather APIs)
and uses a Large Language Model to extract structured disruption events from that
text \textit{before} a commuter travels. Each event carries a severity score
$\sigma$ and a confidence weight $\kappa$. A Temporal Heterogeneous Graph Neural
Network (HGNN) then adjusts those confidence values using city-wide spatial
context. A Bayesian fusion layer converts the signals into per-road disruption
probabilities, which are used to colour-code and rank candidate routes on a
Leaflet.js map.

\vspace{0.5cm}
\noindent
The system supports eight transport modes spanning drive, walk, bike, all five
Kolkata Metro lines, and bus. The API is split into two steps: routes appear on
screen within $\sim$2~s while the slower LLM pipeline ($\sim$30--60~s) fills in
risk scores in the background.

\vspace{0.5cm}

\begin{table}[H]
\centering
\begin{tabular}{ll}
\toprule
\textbf{Aspect} & \textbf{Summary} \\
\midrule
Pilot city          & Kolkata, India \\
Data sources        & 9 (structured + unstructured) \\
Transport modes     & 8 \\
LLM backend         & OpenRouter / Gemini / Ollama (configurable) \\
Graph model         & Temporal HAN -- 3 output heads \\
Probabilistic layer & Bayesian fusion with time-of-day priors \\
\bottomrule
\end{tabular}
\end{table}

\newpage
